\documentclass{ximera}

\input{../../preamble.tex}

\title{Critical Numbers+}

\begin{document}

\begin{abstract}
candidates for change
\end{abstract}
\maketitle








The sign of the derivative tells us where the original function is increasing or decreasing.  Therefore, we are very interested in where the derivative changes sign.

This can happen at three types of numbers. 


\textbf{\textcolor{red!70!darkgray}{$\blacktriangleright$}} The derivative can change signs across a zero. 


\textbf{\textcolor{red!70!darkgray}{$\blacktriangleright$}} The derivative can change signs across a discontiuity. 


\textbf{\textcolor{red!70!darkgray}{$\blacktriangleright$}}  The derivative can change signs across a singularity. 



\begin{center}
The operative word here is \textbf{\textcolor{red!80!black}{CAN}}.
\end{center}




\begin{idea}  \textbf{\textcolor{blue!55!black}{Critical Numbers}}  

Let $f$ be a function with derivative, $f'$. 


A \textbf{critical number} of $f$ is a domain number of $f$ where $f' = 0$ or $f'$ does not exist.


\end{idea}

Functions \textbf{\textcolor{red!70!black}{CAN}} switch behavior at critical numbers. 

They can also switch at singulatrities.






\begin{idea}  \textbf{\textcolor{blue!55!black}{Singularities}}  

Let $f$ be a function. 

A \textbf{singularity} of $f$ is a non-domain number where $f$ has weird behavior.





\textbf{Note:} We have expressed somewhat stable descriptions for discontiuities, but not for singularities.  Calculus will help us with this. The full definition of a singularity will involve the derivative of the derivative of the derivative of the derivative, etc.



\end{idea}

Functions \textbf{\textcolor{red!70!black}{CAN}} switch behavior at singularities. 



\begin{example} Sudden Changes in Slope


Let $K(w)$ be defined by the following graph.



\begin{image}
\begin{tikzpicture} 
  \begin{axis}[
            domain=-10:10, ymax=10, xmax=10, ymin=-10, xmin=-10,
            axis lines =center, xlabel=$w$, ylabel=$y$,
            ytick={-10,-8,-6,-4,-2,2,4,6,8,10},
            xtick={-10,-8,-6,-4,-2,2,4,6,8,10},
            ticklabel style={font=\scriptsize},
            every axis y label/.style={at=(current axis.above origin),anchor=south},
            every axis x label/.style={at=(current axis.right of origin),anchor=west},
            axis on top
          ]
          
          \addplot [line width=2, penColor, smooth, domain=(-4:4),<-] {0.5*x-4};

          \addplot [line width=2, penColor, smooth, domain=(4:7),->] {3*x-14};
        




  \end{axis}
\end{tikzpicture}
\end{image}


This graph makes a sudden change in direction at $(4,-2)$.


On the left side of this point, it appears that the tangent line would be the line $y = \frac{1}{2} \, w - 4$ with a slope of $\frac{1}{2}$.


On the right side of this point, it appears that the tangent line would be the line $y = 3 \, w - 14$ with a slope of $3$.


These are not equal.

$4$ is a critical number of $K$.


However, $K$ is an increasing function. Its behavior did not change at the critical number.



\end{example}


Critical numbers are candidates for a chnage in behavior.





























Hunting down zeros, discontinuities, and singularities have risen to the top of our to-do list.  They are the crux to function behavior. 














\begin{center}
\textbf{\textcolor{green!50!black}{ooooo-=-=-=-ooOoo-=-=-=-ooooo}} \\

more examples can be found by following this link\\ \link[More Examples of the Derivative]{https://ximera.osu.edu/csccmathematics/precalculus2/precalculus2/theDerivative/examples/exampleList}

\end{center}








\end{document}
